% !TEX root =./ECE444_Practice_main.tex
%
% L27 practice set (Null Steering Theory) -- steering vectors, the weight-subtraction
% rule, the PHASER gain/phase table, and the quantization limit on null depth.

\fullwidth{\textbf{Learning Objective 3.9: } Calculate null-steering weights and implement
pattern nulls on the ADALM-PHASER.
\\[4pt]\normalfont\small Show your work. A \textbf{Documentation} line at the end
is required (write \textbf{None} if you did not collaborate).}

\begin{questions}

\question \textbf{Why a null beats gain.} The PHASER views a target on boresight while an
interferer sits at $+22.5\mydeg$. At the aperture the interferer is $10\db$ stronger than the
target. The array is uniformly weighted and steered broadside, so its pattern level at
$+22.5\mydeg$ is the first sidelobe level, $-13\db$ relative to the peak.

\begin{parts}

	\begin{minipage}[t]{0.58\linewidth}
		\part Compute the signal-to-interference ratio at the beamformer output.
		\begin{solution}[1.1in]
		The target sits at the pattern peak and the interferer is attenuated by the
		sidelobe level, so both are referred to the same peak:
		\eq{ \text{SIR} &= 0\db - \lp 10\db - 13\db \rp \\ &= 3\db }
		\end{solution}
	\end{minipage}\hfill%
	\begin{minipage}[t]{0.37\linewidth}
		\raggedleft \vspace{12pt}
		SIR = \ansbox[1.3in]{$3\db$}
	\end{minipage}

	\begin{minipage}[t]{0.58\linewidth}
		\part Null steering replaces the sidelobe at $+22.5\mydeg$ with a notch measured at
		$-21.6\db$ relative to the new main-lobe peak. Recompute the SIR.
		\begin{solution}[1.1in]
		Both quantities are again referred to the pattern peak, so the $1.8\db$ the main lobe
		loses cancels in the ratio:
		\eq{ \text{SIR} &= 0\db - \lp 10\db - 21.6\db \rp \\ &= 11.6\db }
		The null buys $8.6\db$ of margin.
		\end{solution}
	\end{minipage}\hfill%
	\begin{minipage}[t]{0.37\linewidth}
		\raggedleft \vspace{12pt}
		SIR = \ansbox[1.3in]{$11.6\db$}
	\end{minipage}

	\part The same $8.6\db$ of margin could be bought by raising the transmit power. State
	what that costs on a radar, in one or two sentences, and say which fix you would specify.
		\begin{solution}[1.3in]
		A radar echo scales with transmit power, so $8.6\db$ more margin requires about
		$7.2$ times the transmit power, with the amplifier, prime power, cooling, and
		transmitted signature that go with it. The interferer is unaffected by any of that.
		Null steering buys the same margin by recomputing sixteen numbers and costs about
		$2\db$ of main-lobe gain, so the null is the fix to specify whenever the interferer
		angle is known.
		\end{solution}

	\part The interferer is at $+22.5\mydeg$ and the target is on boresight. Explain why the
	array does not reject the interferer simply because it is off axis.
		\begin{solution}[1.1in]
		An array rejects a direction only by the pattern level it presents in that direction.
		Being off axis is not enough; the interferer arrives through whatever sidelobe covers
		$+22.5\mydeg$, and for a uniform eight-element array that sidelobe is only $13\db$
		down.
		\end{solution}

\end{parts}

\newpage

\question \textbf{Steering vectors.} The PHASER array has $N = 8$ elements at
$d = 14\ \text{mm}$, receiving the HB100 tone at $10.525\ \ghz$ so that $d/\lambda = 0.491$
and $kd = 3.085$ rad. The steering vector for arrival angle $\theta$ has entries
$v_n(\theta) = e^{jnkd\sinp{\theta}}$, and the beamformer response is
$y(\theta) = \sum_n w_n v_n(\theta)$.

\begin{parts}

	\part In one or two sentences, state what $\vec{v}(\theta)$ describes and what it does
	\emph{not} depend on.
		\begin{solution}[1.0in]
		The steering vector describes how a plane wave from direction $\theta$ is spread in
		phase across the aperture: element $n$ leads element $0$ by $nkd\sinp{\theta}$. It is
		a property of the arrival direction and the array geometry only, and does not depend
		on the gains and phases loaded into the beamformer.
		\end{solution}

	\begin{minipage}[t]{0.58\linewidth}
		\part Compute the element-to-element phase for a commanded beam at
		$\theta_0 = 30\mydeg$.
		\begin{solution}[1.0in]
		\eq{ \Delta\phi &= kd\sinp{\theta_0} = 3.085\lp 0.500 \rp \\ &= 1.543\ \text{rad}
		= 88.4\mydeg }
		\end{solution}
	\end{minipage}\hfill%
	\begin{minipage}[t]{0.37\linewidth}
		\raggedleft \vspace{12pt}
		$\Delta\phi$ = \ansbox[1.3in]{$88.4\mydeg$}
	\end{minipage}

	\begin{minipage}[t]{0.58\linewidth}
		\part Write the weight $w_3$ (the fourth element, $n = 3$) that steers the beam to
		$30\mydeg$, as a magnitude and an angle.
		\begin{solution}[1.1in]
		Beam steering uses the conjugate steering vector, $w_n = e^{-jnkd\sinp{\theta_0}}$,
		so every element runs at full gain and element $n$ is retarded by $n\Delta\phi$:
		\eq{ w_3 = e^{-j3(88.4\mydeg)} = 1\ \angle -265.2\mydeg = 1\ \angle +94.8\mydeg }
		\end{solution}
	\end{minipage}\hfill%
	\begin{minipage}[t]{0.37\linewidth}
		\raggedleft \vspace{12pt}
		$w_3$ = \ansbox[1.3in]{$1\ \angle +94.8\mydeg$}
	\end{minipage}

	\begin{minipage}[t]{0.58\linewidth}
		\part With those weights loaded, evaluate $y(30\mydeg)$ and express the result in dB
		relative to one element.
		\begin{solution}[1.1in]
		Each term becomes $e^{-jnkd\sinp{\theta_0}}e^{+jnkd\sinp{\theta_0}} = 1$, so the eight
		phasors add in phase:
		\eq{ y(30\mydeg) &= \sum_{n=0}^{7} 1 = 8 \\
		20\mylog{8} &= 18.1\db }
		\end{solution}
	\end{minipage}\hfill%
	\begin{minipage}[t]{0.37\linewidth}
		\raggedleft \vspace{12pt}
		$y$ = \ansbox[1.3in]{$8$ ($18.1\db$)}
	\end{minipage}

\end{parts}

\newpage

\question \textbf{Weight subtraction by hand.} A four-element array with $d = \lambda/2$
(so $kd = 180\mydeg$) is steered broadside, giving $\vec{w}_\text{d} = [1, 1, 1, 1]$. An
interferer appears at $\theta_1 = 45\mydeg$. Use the weight-subtraction rule
$\vec{w} = \vec{w}_\text{d} - r_\text{n}\vec{w}_\text{n}$ with
$r_\text{n} = \vec{w}_\text{n}^{H}\vec{w}_\text{d} / \vec{w}_\text{n}^{H}\vec{w}_\text{n}$.

\begin{parts}

	\begin{minipage}[t]{0.58\linewidth}
		\part Compute the element-to-element phase $\psi_1 = kd\sinp{\theta_1}$ at the
		interferer angle.
		\begin{solution}[1.0in]
		\eq{ \psi_1 &= 180\mydeg \times \sinp{45\mydeg} \\ &= 180\mydeg (0.7071)
		= 127.3\mydeg }
		\end{solution}
	\end{minipage}\hfill%
	\begin{minipage}[t]{0.37\linewidth}
		\raggedleft \vspace{12pt}
		$\psi_1$ = \ansbox[1.3in]{$127.3\mydeg$}
	\end{minipage}

	\begin{minipage}[t]{0.58\linewidth}
		\part Evaluate $r_\text{n}$. For a broadside beam it reduces to the normalized
		uniform array factor at $\theta_1$, $r_\text{n} = e^{j(N-1)\psi_1/2}\ \sinp{N\psi_1/2}/\lp N\sinp{\psi_1/2} \rp$.
		\begin{solution}[1.2in]
		With $N = 4$ and $\psi_1 = 127.3\mydeg$:
		\eq{ \frac{\sinp{2\psi_1}}{4\sinp{\psi_1/2}} &= \frac{\sinp{254.6\mydeg}}
		{4\sinp{63.6\mydeg}} = \frac{-0.964}{3.585} = -0.269 \\
		r_\text{n} &= (-0.269)\ e^{j190.9\mydeg} = 0.269\ \angle 10.9\mydeg }
		The magnitude $0.269$ is $-11.4\db$, which is the uniform four-element pattern level
		at $45\mydeg$.
		\end{solution}
	\end{minipage}\hfill%
	\begin{minipage}[t]{0.37\linewidth}
		\raggedleft \vspace{12pt}
		$r_\text{n}$ = \ansbox[1.3in]{$0.269\ \angle 10.9\mydeg$}
	\end{minipage}

	\part Compute the four weights $w_n = 1 - r_\text{n}e^{-jn\psi_1}$ and convert them to
	PHASER-style settings: element gains as percentages of the largest, and phases in degrees.
		\begin{solution}[1.9in]
		The subtracted phasor has magnitude $0.269$ and angle $10.9\mydeg - n(127.3\mydeg)$:
		\eq{ w_0 &= 1 - 0.269\angle 10.9\mydeg = 0.736 - j0.051 = 0.738\ \angle -4.0\mydeg \\
		w_1 &= 1 - 0.269\angle -116.4\mydeg = 1.119 + j0.241 = 1.145\ \angle +12.2\mydeg \\
		w_2 &= 1 - 0.269\angle -243.7\mydeg = 1.119 - j0.241 = 1.145\ \angle -12.2\mydeg \\
		w_3 &= 1 - 0.269\angle -371.0\mydeg = 0.736 + j0.051 = 0.738\ \angle +4.0\mydeg }
		Scaling the largest magnitude to 100\%, the settings are
		\eq{ \text{gains} &= 64, 100, 100, 64\ \text{percent} \\
		\text{phases} &= -4.0\mydeg,\ +12.2\mydeg,\ -12.2\mydeg,\ +4.0\mydeg }
		The gains are symmetric and the phases antisymmetric, as they must be
		for a single null placed off a broadside beam.
		\end{solution}

	\begin{minipage}[t]{0.58\linewidth}
		\part Find the loss in the look direction after the weights are rescaled so the
		largest element sits at 100\% gain.
		\begin{solution}[1.2in]
		At broadside every steering-vector entry is 1, so the response is the plain sum of
		the weights, and the rescaling divides by the largest magnitude:
		\eq{ \sum_n w_n &= 3.711, \qquad \max_n \vert w_n \vert = 1.145 \\
		\text{loss} &= 20\mylog{\frac{3.711}{1.145 \times 4}} = -1.8\db }
		\end{solution}
	\end{minipage}\hfill%
	\begin{minipage}[t]{0.37\linewidth}
		\raggedleft \vspace{12pt}
		loss = \ansbox[1.3in]{$1.8\db$}
	\end{minipage}

\end{parts}

\newpage

\question \textbf{What the null costs.} Return to the eight-element PHASER array
($kd = 3.085$ rad, broadside beam). Before the weights are rescaled, the response in the
look direction falls from $N$ to $N\lp 1 - \vert r_\text{n} \vert^2 \rp$.

\begin{parts}

	\begin{minipage}[t]{0.58\linewidth}
		\part An interferer sits at $\theta_1 = 10\mydeg$, where the uniform pattern level is
		$-8.0\db$. Compute the loss in the look direction before rescaling.
		\begin{solution}[1.1in]
		A level of $-8.0\db$ means $\vert r_\text{n} \vert = 10^{-8.0/20} = 0.398$:
		\eq{ \text{loss} &= 20\mylog{1 - (0.398)^2} \\ &= 20\mylog{0.842} = -1.5\db }
		\end{solution}
	\end{minipage}\hfill%
	\begin{minipage}[t]{0.37\linewidth}
		\raggedleft \vspace{12pt}
		loss = \ansbox[1.3in]{$1.5\db$}
	\end{minipage}

	\begin{minipage}[t]{0.58\linewidth}
		\part A second interferer sits at $\theta_1 = 30\mydeg$, where the uniform pattern
		level is $-33.9\db$. Compute the loss and comment on the result in one sentence.
		\begin{solution}[1.2in]
		Here $\vert r_\text{n} \vert = 10^{-33.9/20} = 0.020$:
		\eq{ \text{loss} = 20\mylog{1 - (0.020)^2} = -0.004\db }
		The interferer is already sitting in one of the array's natural pattern nulls, so
		the rule subtracts almost nothing and the null is free.
		\end{solution}
	\end{minipage}\hfill%
	\begin{minipage}[t]{0.37\linewidth}
		\raggedleft \vspace{12pt}
		loss = \ansbox[1.3in]{$\approx 0\db$}
	\end{minipage}

	\part The PHASER's half-power beamwidth is $13.1\mydeg$. A student asks the array to null
	an interferer at $\theta_1 = 3\mydeg$, where $\vert r_\text{n} \vert = 0.93$. Explain what
	the resulting pattern looks like and what the student should do instead.
		\begin{solution}[1.3in]
		With $\vert r_\text{n} \vert = 0.93$ the rule subtracts almost the entire beam. The
		eight weights very nearly cancel, and after rescaling what remains is a split pattern
		with a hole where the target used to be; the look-direction gain falls by more than
		$12\db$ relative to the uniform array. An interferer inside the half-power beamwidth cannot be separated from the
		target by weighting, so the student should steer the beam to a different angle or
		wait for the geometry to change.
		\end{solution}

	\part Explain, in one or two sentences, why $\vert r_\text{n} \vert$ is worth computing
	before any weights are calculated.
		\begin{solution}[1.0in]
		$\vert r_\text{n} \vert$ is the level of the pattern you already have at the null
		angle, and it sets the main-lobe loss. Reading it off the existing pattern tells you
		whether the null is cheap or ruinous before any weight arithmetic is done.
		\end{solution}

\end{parts}

\newpage

\question \textbf{How deep can the null be?} Each ADAR1000 phase setting lands on a
$2.8125\mydeg$ grid and each gain on a 1\% grid, so each weight carries an independent error.
The errors add in RMS at the null angle, giving a residual response
$\vert y(\theta_1)\vert / N \approx \epsilon_\text{rms}/\sqrt{N}$ with
$\epsilon_\text{rms} = \sqrt{\sigma_\phi^2 + \sigma_a^2}$ and
$\sigma_\phi = \text{LSB}/\sqrt{12}$ in radians.

\begin{parts}

	\begin{minipage}[t]{0.58\linewidth}
		\part Compute the quantization-limited null depth for the eight-element PHASER with
		its 7-bit phase shifters.
		\begin{solution}[1.3in]
		\eq{ \sigma_\phi &= \frac{2.8125\mydeg}{\sqrt{12}} = 0.812\mydeg = 0.0142\ \text{rad}
		\\ \sigma_a &= \frac{0.01}{\sqrt{12}} = 0.0029 \\
		\epsilon_\text{rms} &= \sqrt{0.0142^2 + 0.0029^2} = 0.0145 \\
		\text{depth} &= 20\mylog{\frac{0.0145}{\sqrt{8}}} = -45.8\db }
		\end{solution}
	\end{minipage}\hfill%
	\begin{minipage}[t]{0.37\linewidth}
		\raggedleft \vspace{12pt}
		depth = \ansbox[1.3in]{$-45.8\db$}
	\end{minipage}

	\begin{minipage}[t]{0.58\linewidth}
		\part Repeat for a 3-bit phase shifter, ignoring the gain error. Compare your answer
		with the $-6B$ rule of thumb from L26.
		\begin{solution}[1.3in]
		A 3-bit part has $\text{LSB} = 360\mydeg/8 = 45\mydeg$:
		\eq{ \sigma_\phi &= \frac{45\mydeg}{\sqrt{12}} = 13.0\mydeg = 0.227\ \text{rad} \\
		\text{depth} &= 20\mylog{\frac{0.227}{\sqrt{8}}} = -21.9\db }
		The $-6B$ rule predicts $-18\db$ for three bits, so the null floor lands a few
		decibels deeper because the error is averaged over eight elements.
		\end{solution}
	\end{minipage}\hfill%
	\begin{minipage}[t]{0.37\linewidth}
		\raggedleft \vspace{12pt}
		depth = \ansbox[1.3in]{$-21.9\db$}
	\end{minipage}

	\part The PHASER carries 7-bit phase shifters, yet the lab sweep measures a notch of only
	$-21.6\db$ relative to the peak. Explain the discrepancy, given that the sweep's noise
	floor sits $23\db$ below the uniform-taper peak and the null-steered main lobe sits
	$2\db$ below that same reference.
		\begin{solution}[1.3in]
		The weights are good to about $-46\db$, so quantization is not what the plot is
		showing. The sweep cannot display anything below its own noise floor, which is
		$23\db$ below the uniform peak. Measured against the null-steered main lobe, which is
		$2\db$ lower, that floor is about $21\db$ down, and $-21.6\db$ is the floor rather
		than the null. Deeper nulls exist in the weights but cannot be seen without lowering
		the noise floor.
		\end{solution}

	\part State the maximum number of independent nulls an eight-element array can place and
	explain the counting in one sentence.
		\begin{solution}[1.0in]
		Seven. The array has eight complex weights, one of which is spent holding the main
		beam in the look direction, leaving $N - 1 = 7$ degrees of freedom for nulls.
		\end{solution}

\end{parts}

\vspace{1.5em}
\noindent\textbf{Documentation:} \rule[-2pt]{0.6\linewidth}{0.4pt}

\end{questions}
